\documentclass[]{article}\usepackage{amsmath}\usepackage{multirow}\usepackage{natbib}\usepackage{graphicx} \begin{document}The identity of dark matter constitutes one of the most significant unsolved problems in fundamental physics. The standard cosmological paradigm, Lambda Cold Dark Matter ($\Lambda$CDM), provides an excellent description of the universe on large scales, from the anisotropies in the cosmic microwave background to the large-scale distribution of galaxies. However, the simple, non-interacting nature of cold dark matter is less rigorously tested on small scales, leaving a vast discovery space for new microphysical phenomena. Among the most compelling theoretical possibilities is the existence of a non-gravitational interaction between dark matter and baryons. If present, such an interaction would have left a distinct and potentially observable imprint on the formation of the very first cosmic structures.In the early universe, elastic scattering between dark matter and baryons would facilitate an exchange of momentum and heat, coupling the two fluids \citep{dvorkin2013constrainingdarkmatterbaryonscattering,das2025darksecretsbaryonsilluminating}. This coupling effectively acts as a pressure that counteracts gravitational collapse below a characteristic length scale, which is directly determined by the dark matter-baryon scattering cross-section \citep{slatyer2018earlyuniverseconstraintsdarkmatterbaryon}. The primary cosmological consequence is a suppression of the matter power spectrum at high wavenumbers \citep{zhang2024weaklensingconstraintsdark,das2025darksecretsbaryonsilluminating}, which translates into a cutoff in the halo mass function \citep{dvorkin2013constrainingdarkmatterbaryonscattering}. This cutoff prevents the formation of the lowest-mass halos, thereby providing a direct link between the microphysics of dark matter and the abundance of small-scale structures \citep{dvorkin2013constrainingdarkmatterbaryonscattering}. Searching for evidence of such a cutoff is therefore a powerful method for probing the fundamental properties of dark matter.The 21 cm forest offers a uniquely direct window into this small-scale regime \citep{ciardi201521cmforestska,shimabukuro2026topologicalsignaturesheatingdark}. Arising from the absorption of radio waves from distant sources by intervening clouds of neutral hydrogen, the 21 cm forest traces the distribution of gas within the high-redshift intergalactic medium and inside the first gravitationally bound objects, known as minihalos \citep{ciardi201521cmforestska,shimabukuro2026topologicalsignaturesheatingdark}. These minihalos are precisely the structures whose formation is most affected by a suppression of small-scale power \citep{ciardi201521cmforestska,shimabukuro2025analyzing21cmforestwavelet}. A key challenge, however, is that the statistics of 21 cm absorption lines are also sensitive to astrophysical effects, such as the intensity of the ambient ionizing background or the thermal state of the gas \citep{ciardi201521cmforestska,šoltinský2021detectabilitystrong21centimetre}, creating potential degeneracies. The central problem is to develop an observational strategy that can robustly distinguish the signature of dark matter interactions from these astrophysical uncertainties \citep{sun2025deeplearningdrivenlikelihoodfreeparameter,šoltinský2025prospectsstatisticaldetection21cm,zhao202621cmforestonedimensional}.In this paper, we quantify the optimal strategy for detecting the signature of dark matter-baryon scattering with the 21 cm forest \citep{hou2025effectsdarkmatterannihilation}. Using a semi-analytic framework to model the statistics of 21 cm absorbers, we demonstrate that the suppression of low-mass minihalos imprints a distinct feature on the shape of the optical depth distribution. This feature, a characteristic deficit of low optical depth systems, provides a powerful discriminant that allows the effects of dark matter microphysics to be disentangled from those of astrophysical parameters \citep{hou2025effectsdarkmatterannihilation}. We perform a Fisher forecast that incorporates a realistic, redshift-dependent evolution for the population of background radio sources to map the sensitivity of this probe across cosmic time. Our analysis identifies an optimal observational window at redshifts $z \sim 8-10$, where the combination of intrinsic physical sensitivity and statistical power is maximized. This work establishes a clear observational path for using next-generation radio observatories to place leading constraints on the nature of dark matter \citep{gluscevic2018constraintsscatteringkevtevdark,straight2026cmbconstraintsdarkmatterproton}.\end{document}                